Sign language datasets can be categorized into isolated or word-level, where each video matches a specific sign, and continuous or sentence-level, where videos contain many signs that correspond to an entire sentence. As mentioned before, while isolated datasets can be used for recognition they cannot be used for translation \cite{bragg2019sign}. Since the main goal of this work is to present a dataset for SLT, only CSL datasets will be considered.

One of the most relevant SLT datasets is RWTH-Phoenix-Weather 2014 T \cite{camgoz2018neural}. It contains videos of German Sign Language (GSL) extracted from German public TV weather forecasts. This dataset is used today as the main benchmark for SLT and, having a vocabulary of over 1000 signs, is considered the only resource for large-scale continuous sign language worldwide \cite{koller2020quantitative}. As the present work, it has the peculiarity of having been recorded in real life conditions, which may result in a more challenging dataset than if it was laboratory-made. Most of the other available SLT datasets have been recorded in laboratory conditions and contain a set of frequent or relevant sentences for its respective language. This is the case of the SIGNUM dataset \cite{von2008significance} of GSL, the Chinese Sign Language dataset (CSL) \cite{huang2018video}, the Greek Sign Language (GSL) dataset \cite{adaloglou2020comprehensive} and the KETI dataset \cite{ko2019neural} of Korean Sign Language.

Table \ref{tab:bd_comparison}  shows the main features of the mentioned datasets alongside \dbname. Many datasets also provide information about glosses translation (an intermediate translation between sign language and spoken language where each gloss corresponds to one sign) so we differentiated between gloss (gl) and word (w) level. There is no information about glosses in \dbname because it only contains sentence-level Spanish translation. There is no gloss information about CLS dataset either, as there is no gloss notation in CSL, signs map directly to words.

\begin{table}
\caption{Main features of the studied datasets.}
\label{tab:bd_comparison}
\begin{center}
    \begin{tabular}{ |l|c|c|c|c|c|c| }
        \hline
        & PHOENIX* & SIGNUM & CSL & GSL & KETI & \dbname\\ 
        \hline
        language & German & German & Chinese & Greek & Korean & Spanish\\
        sign language & GSL & GSL & CSL & GSL & KLS & LSA\\
        real life & \textbf{Yes} & No & No & No & No & \textbf{Yes}\\
        signers & 9 & 25 & 50 & 7 & 14 & \textbf{103}\\
        duration [h] & 10.71 & 55.3 & \textbf{100+} & 9.51 & 28 & 21.78\\ 
        \# samples & 7096 & \textbf{33,210} & 25,000 & 10,295 & 14,672 & 14,880\\ 
        \hline
        \# unique sentences & 5672 & 780 & 100 & 331 & 105 & \textbf{14,254}\\
        \% unique sentences & 79.93\% & 2.35\% & 0.4\% & 3.21\% & 0.71\% & \textbf{95.79\%}\\
        vocab. size (w) & 2887 & N/A & 178 & N/A & 419 & \textbf{14,239}\\ 
        \# singletons (w) & 1077 & 0 & 0 & 0 & 0 & \textbf{7150}\\
        \% singletons (w) & 37.3\% & 0\% & 0\% & 0\% & 0\% & \textbf{50.21\%}\\
        vocab. size (gl) & \textbf{1066} & 450 & - & 310 & 524 & -\\
        \# singletons (gl) & \textbf{337} & 0 & - & 0 & 0 & -\\
        \# singletons (gl) & \textbf{31.61\%} & 0\% & - & 0\% & 0\% & -\\
        \hline
        resolution & 210x260 & 776x578 & \textbf{1920x1080} & 848x480 & \textbf{1920x1080} & \textbf{1920x1080}\\
        fps & 25 & \textbf{30} & \textbf{30} & \textbf{30} & \textbf{30} & \textbf{30}\\
        \hline
    \end{tabular}
*Data was not available for the whole PHOENIX dataset, so the table show its train set statistics.
\end{center}
\end{table}

It is interesting to notice that there is a significant difference between datasets generated in real life conditions and those generated in laboratory conditions regarding the amount of unique sentences, the vocabulary size and the amount of singletons (tokens that appear only once in the training set). As laboratory-made dataset first choose a number of sentences and then record them many times, the set of unique sentences is a small fraction of the total amount of samples (between 0 and 5 percent). On the other hand, as it is not as common to find an exact sentence repeated many times in realistic environments, datasets generated in real life conditions present a much higher rate of unique sentences: 79.93\% for the PHOENIX dataset and 95.79\% for LSA-T. This rate was expected to be lower in PHOENIX than in LSA-T as the videos in PHOENIX were taken only from the weather forecast, while videos in LSA-T feature a wide range of topics. This directly impacts the size of the vocabulary and the amount of unique words in each dataset. While laboratory-made datasets have a small vocabulary with no singletons (all sentences appear multiple times), datasets collected in real-life environments have large vocabulary sizes with a significant amount of unique words (37.3\% in PHOENIX and 50.21\% in LSA-T). As a consequence, datasets recorded in real-life conditions, and particularly the one presented in this work, are highly challenging for translation models.

\subsection{LSA}

LSA is the sign language used by the deaf community in Argentina. In 2018, 0.02\% of the Argentinian population older than 6 years old presented a hearing difficulty, half of them being deaf \cite{instituto2018estudio}. Counting teachers, relatives and friends of deaf people, there are 2 million LSA users \cite{url:lsa}.

Regarding LSA datasets, LSA64 \cite{ronchetti2016sign}\cite{ronchetti2016lsa64} is, to the best of our knowledge, the only other dataset available. It contains 3200 videos of 64 isolated signs recorded by a team of researchers from the Universidad Nacional de La Plata. As it has been generated under laboratory conditions, the background scenarios are similar. Although it has been the only source of data for training LSA sign recognition models, it cannot be used for translation.

\subsection{SLT}

Regarding SLT, works that use keypoints data as part of the input for a translation model can be found in \cite{borg2020phonologically} and  \cite{zhou2020spatial}. In those cases keypoints data were given to the model as extra information together with the corresponding video. Furthermore, the models presented in \cite{ko2019neural} and \cite{kim2022keypoint} have a similar approach to the one presented in this work. Also, they only use keypoints information as input for a neural network model. In both cases the authors train a recurrent network on the KETI dataset. Particularly, in \cite{kim2022keypoint} no gloss information is employed, neither as input or output.
